\begin{center}
    \textbf{\large Вариант 3}
\end{center}

\samerowans{\begin{problem}{1}
Вычислите: $5104-4354 : (252-14\cdot17)+108$
\end{problem}}

\vspace{5pt}

\samerowans{\begin{problem}{2}
Решите уравнение: $(360 : a - 2 \cdot 15) \cdot 5 = 150$
\end{problem}}

\vspace{5pt}

\samerowans{\begin{problem}{3}
Вычислите: $6\text{ т } 84\text{ кг } - 45\text{ ц } + 126\text{ кг}$. Ответ запишите в т, ц, кг.
\end{problem}}

\samerowans{\begin{problem}{4}
Сложили пять одинаковых чисел, к результату прибавили 10, затем полученную сумму разделили на 8 и получили 15. Какие числа складывали?
\end{problem}}

\vspace{5pt}

\samerowans{\begin{problem}{5}
Катя два часа делала домашнее задание. Две трети часа она делала математику, одну четверть часа — русский язык. Все оставшееся время она учила английский. Сколько минут Катя учила английский?
\end{problem}}

\vspace{5pt}

\begin{problem}{6} Из двух пунктов А и В, расстояние между которыми 60 км, одновременно навстречу друг другу выехали два велосипедиста. Скорость велосипедиста, выехавшего из п А, на $5\text{ км/ч}$ больше скорости второго. За два часа им удалось встретиться и разъехаться на расстояние 30 км. На каком расстоянии от п А будет второй велосипедист через 4 часа 15 минут после своего выезда из В?
\end{problem}

\vspace{5pt}

\noindent\grid{36}{19}

\vspace{5pt}

\samerowans{\begin{problem}{7}
Первый станок может сделать 420 деталей за 6 часов, второй справится с этой задачей в 2 раза быстрее. Сколько часов совместной работы потребуется двум станкам, чтобы сделать 630 деталей?
\end{problem}}

\newpage

\begin{problem}{8} Из прямоугольника с периметром 22 дм вырезали квадрат со стороной 4 см. Найди площадь оставшейся фигуры, если длина прямоугольника 700 мм.
\end{problem}

\smallskip

\noindent\grid{36}{15}

\vspace{5pt}

\samerowans{\begin{problem}{9}
Таня заплатила за покупку 1 м 40 см шелковой ленты 420 рублей. Надя купила ленты на 50 см больше, чем Таня. Сколько рублей заплатила Надя?
\end{problem}}

\vspace{5pt}

\samerowans{\begin{problem}{10}
Коля и Петя сыграли партию в шахматы, которая закончилась победой одного из них. «Петя выиграл», — сказал Коля. «Я победил», — сказал Петя. Кто из них выиграл, если известно, что хотя бы кто-то из них сказал неправду?
\end{problem}}

\vspace{5pt}

\begin{problem}{11} Через 8 лет мальчик будет вдвое старше, чем он был 3 года назад. А девочка будет через 4 года втрое старше, чем 8 лет назад. Кто старше: мальчик или девочка?
\end{problem}

\smallskip

\noindent\grid{36}{20}

\newpage

\begin{center}
\large{Вариант 3}
\end{center}

\begin{enumerate}
    \item \textbf{Вычислите:} $5104 - 4354 : (252 - 14 \cdot 17) + 108$ \\
    \textit{Решение:}
    \begin{itemize}
        \item $14 \cdot 17 = 238$
        \item $252 - 238 = 14$
        \item $4354 : 14 = 311$
        \item $5104 - 311 + 108 = 4901$
    \end{itemize}
    \textbf{Ответ: 4901}

    \item \textbf{Решите уравнение:} $(360:a - 2 \cdot 15) \cdot 5 = 150$ \\
    \textit{Решение:}
    \begin{itemize}
        \item $360:a - 30 = 150 : 5$
        \item $360:a - 30 = 30$
        \item $360:a = 60 \implies a = 6$
    \end{itemize}
    \textbf{Ответ: 6}

    \item \textbf{Вычислите:} 6 т 84 кг $- 45$ ц $+ 126$ кг \\
    \textit{Решение:}
    \begin{itemize}
        \item $6084 \text{ кг} - 4500 \text{ кг} + 126 \text{ кг} = 1710 \text{ кг}$
        \item $1710 \text{ кг} = 1 \text{ т } 7 \text{ ц } 10 \text{ кг}$
    \end{itemize}
    \textbf{Ответ: 1 т 7 ц 10 кг}

    \item \textbf{Задача (пять одинаковых чисел):} \\
    \textit{Решение:} Пусть $x$ — искомое число. \\
    $(5x + 10) : 8 = 15 \implies 5x + 10 = 120 \implies 5x = 110 \implies x = 22$ \\
    \textbf{Ответ: 22}

    \item \textbf{Домашнее задание Кати:} \\
    \textit{Решение:} Общее время 2 часа = 120 мин.
    \begin{itemize}
        \item Математика: $\frac{2}{3} \cdot 60 = 40$ мин
        \item Русский: $\frac{1}{4} \cdot 60 = 15$ мин
        \item Английский: $120 - (40 + 15) = 65$ мин
    \end{itemize}
    \textbf{Ответ: 65 минут}

    \item \textbf{Велосипедисты:} \\
    \textit{Решение:}
    \begin{itemize}
        \item Суммарный путь: $60 + 30 = 90$ км. Скорость сближения: $90 / 2 = 45$ км/ч.
        \item $v_1 + v_2 = 45$ и $v_1 = v_2 + 5 \implies v_2 = 20, v_1 = 25$ км/ч.
        \item Путь второго из B за 1 ч 15 мин: $20 \cdot 1.25 = 25$ км.
        \item Расстояние от A: $60 - 25 = 35$ км.
    \end{itemize}
    \textbf{Ответ: 35 км}

    \item \textbf{Станки:} \\
    \textit{Решение:}
    \begin{itemize}
        \item Производительность 1: $420 / 6 = 70$ дет/ч. Производительность 2: $420 / 3 = 140$ дет/ч.
        \item Время совместной работы: $630 / (70 + 140) = 3$ часа.
    \end{itemize}
    \textbf{Ответ: 3 часа}

    \item \textbf{Площадь фигуры:} \\
    \textit{Решение:} $P = 220$ см, длина $a = 70$ см.
    \begin{itemize}
        \item Ширина $b = (220 - 2 \cdot 70) / 2 = 40$ см.
        \item $S = (70 \cdot 40) - (4 \cdot 4) = 2800 - 16 = 2784 \text{ см}^2$.
    \end{itemize}
    \textbf{Ответ: 2784 см$^2$}

    \item \textbf{Шелковая лента:} \\
    \textit{Решение:} 140 см стоят 420 руб $\implies$ 1 см стоит 3 руб. \\
    Надя купила $140 + 50 = 190$ см. Стоимость: $190 \cdot 3 = 570$ руб. \\
    \textbf{Ответ: 570 рублей}

    \item \textbf{Шахматы (Логика):} \\
    \textit{Решение:} Оба утверждения означали победу Пети. Так как по условию один из них солгал, победа Пети невозможна. Выиграл Коля. \\
    \textbf{Ответ: Коля}

    \item \textbf{Возраст:} \\
    \textit{Решение:}
    \begin{itemize}
        \item Мальчик: $x + 8 = 2(x - 3) \implies x = 14$ лет.
        \item Девочка: $y + 4 = 3(y - 8) \implies y = 14$ лет.
    \end{itemize}
    \textbf{Ответ: одинаково (по 14 лет)}
\end{enumerate}